\section{Rekam Medis Elektronik}

\subsection{Pengertian}
Berdasarkan Permenkes nomor 269 Tahun 2008, Rekam medis adalah berkas yang berisikan catatan dan dokumen tentang identitas pasien, pemeriksaan, pengobatan, tindakan dan pelayanan lain yang telah diberikan kepada pasien. Rekam medis dapat dipisah menjadi dua tipe yaitu rekam medis pasien rawat jalan dan rawat inap dengan perbedaanya pada tambahan informasi persetujuan tindakan medis pada rawat inap \parencite{1}.

Rekam medis sendiri dibuat untuk memenuhi 2 jenis tujuan yaitu tujuan utama dan tujuan sekunder. Tujuan utama dari rekam medis terdiri atas lima kepentingan yaitu kepentingan pasien, pelayanan pasien, manajemen kesehatan, menunjang pelayanan, dan pembiayaan. Sementara itu, tujuan sekunder dari rekam medis lebih ditujukan kepada hal-hal terkait dengan lingkungan pelayanan pasien seperti administrasi, hukum, keuangan, penelitian, pendidikan dan dokumentasi \parencite{1}.

\begin{itemize}
    \item Pengobatan Pasien.
    \item Peningkatan Kualitas Pelayanan.
    \item Pendidikan dan Penelitian.
    \item Pembiayaan.
    \item Statistik Kesehatan.
    \item Pembuktian Masalah Hukum,Disiplin, dan Etik.
\end{itemize}

\subsection{Komponen dan Struktur Rekam Medis Elektronik}


Sistem informasi rumah sakit sering kali menghadapi masalah kualitas dalam pencatatan jejak audit (\textit{audit trails}). Analisis terhadap sistem informasi rumah sakit menunjukkan bahwa mekanisme audit yang ada sering kali tidak memadai untuk merekonstruksi insiden keamanan secara utuh atau menjamin integritas data jangka panjang \cite{6}.

Selain itu, integrasi teknologi baru seperti \textit{blockchain} dalam layanan kesehatan telah diidentifikasi memiliki potensi keuntungan sekaligus ancaman baru. Meskipun menjanjikan desentralisasi, penerapannya harus mempertimbangkan interoperabilitas dan skalabilitas dalam lingkungan medis yang kompleks \cite{1}.

\subsection{Standar dan Regulasi Rekam Medis Elektronik di Indonesia}
Di Indonesia, penyelenggaraan RME diatur secara ketat oleh Peraturan Menteri Kesehatan (Permenkes) No. 24 Tahun 2022 \cite{1}. Peraturan ini mewajibkan seluruh fasilitas pelayanan kesehatan untuk beralih ke RME dan terhubung dengan platform SATUSEHAT, yang bertujuan untuk menciptakan interoperabilitas data kesehatan nasional. Untuk mencapai interoperabilitas ini, standar internasional menjadi acuan penting. Standar seperti ISO 18308 (\textit{Health informatics — Requirements for an electronic health record architecture}) \cite{12} dan HL7 FHIR (\textit{Fast Healthcare Interoperability Resources}) \cite{11} menjadi acuan global untuk memastikan data dapat dipertukarkan antar sistem secara konsisten, aman, dan semantik.

\subsection{Tantangan Keamanan dan Privasi EMR}
Data RME bersifat sangat rahasia dan bernilai tinggi, menjadikannya target utama serangan siber \cite{17}. Tantangan utama yang dihadapi meliputi risiko \textit{data breach} (pencurian data), akses tidak sah baik oleh pihak internal maupun eksternal, serta ancaman terhadap integritas data. Untuk mengatasi ini, sistem RME harus mengimplementasikan kontrol keamanan berlapis, termasuk enkripsi data (baik saat transit maupun saat disimpan), mekanisme audit log yang komprehensif untuk melacak semua aktivitas, serta jaminan \textit{non-repudiation} (kenirsangkalan) untuk memastikan setiap tindakan dapat diatribusikan secara unik dan tidak dapat disangkal oleh pelakunya.